% Printed 401
\par\clearpage
\noindent\null\hfill{\small\textbf{ON THE HANOVER TRIANGULATION.}}\hfill\makebox[0pt][r]{\fontsize{13}{16}\selectfont 401}\par
\vspace*{4\baselineskip}
\begin{center}
{\large III.}\par\medskip
{\LARGE\textls[150]{POSTHUMOUS PAPERS}}\par\smallskip
{\large ON THE HANOVER TRIANGULATION.}
\end{center}
\vspace{2\baselineskip}
\begin{center}
[1.]\par\smallskip
{\large [Plan and Opening of the Work on the Trigonometric Surveys\\
in Hanover.]}\par\medskip
\textls[120]{Plan of the Work.}\par\smallskip
Eight Sections.
\end{center}

\noindent
\begin{tabular*}{\linewidth}{@{}r@{\enspace}l@{\extracolsep{\fill}}r@{\hspace{3.6em}}}
1) & On the arrangements and methods used in the surveys & \\
   & in general & \makebox[0pt][r]{about\enspace}40\makebox[0pt][l]{\enspace pages}\\
2) & Measurement of the angles of the principal triangles & 60\\
3) & Minor reductions of the angle measurements & 12\\
4) & Adjustment of the measurements & 24\\
5) & First method of computing the triangles, as represented & \\
   & on the spheroid & 36\\
6) & Second method of computation, as represented in the plane & 60\\
7) & Secondary points & 24\\
8) & Height measurements & 20\\
\multicolumn{3}{@{}r@{}}{A total of about 276 pages}\\
\multicolumn{3}{@{}r@{}}{or \GaussIXMathCopy{34 1/2}{\(34\tfrac{1}{2}\)} sheets.}
\end{tabular*}\par\medskip

\noindent\hspace*{\parindent}The first section is to deal with the following subjects: immediate purpose
of the surveys; distinctive features of the terrain; target points; observation stations;
instruments; method of observation; baseline; provisional computation of the points.
\par\smallskip
\noindent\hspace*{.031\linewidth}{\footnotesize IX.}\hfill{\small 51}\hspace*{.128\linewidth}\par

% Printed 402
\clearpage
\noindent\makebox[0pt][l]{\fontsize{13}{16}\selectfont 402}\hfill{\small\textbf{POSTHUMOUS PAPERS.}}\hfill\null\par
\vspace{2.5\baselineskip}
\begin{center}
{\large\textls[120]{The Trigonometric Surveys in the Kingdom of Hanover.}}\par\medskip
\rule{.12\linewidth}{.4pt}\par\medskip
{\large\textls[120]{Introduction.}}
\end{center}

\par\noindent\hspace*{\parindent}The first occasion for the trigonometric surveys I carried out in the Kingdom
of Hanover in the years 1821---1825 was provided by the arc measurement
undertaken in the Danish states by Councillor of State \textsc{Schumacher}.
While the northern tip of Jutland was bound to be its natural endpoint,
it could be extended southward, reaching its limit only at the Mediterranean Sea. Such an extension
would first have to cross the Kingdom of Hanover over a stretch of more
than two degrees of latitude. At the kingdom's southern border, the Göttingen
Observatory afforded the opportunity for the most precise and convenient
astronomical determinations. Our enlightened government could hardly fail
to recognize the scientific importance of such an undertaking and the firm
foundation it was bound to provide for the geography of the kingdom
and for future, more extensive trigonometric surveys. I was therefore
commissioned in 1820 to carry out this continuation of the Danish arc measurement
through the Kingdom of Hanover.

\par\noindent\hspace*{\parindent}The trigonometric part of this work began at the southern end in the
summer of 1821 and was completed in the summer of 1823 with the measurements
on Michaelisthurm in Hamburg.

\par\noindent\hspace*{\parindent}In the meantime a trigonometric survey of the neighboring Electorate
of Hesse had been begun under the direction of Professor \textsc{Gerling}. The
principal triangles were to be measured with exceptionally good equipment
and with all attainable accuracy. Connecting them with the Hanoverian
triangles was therefore all the more important, since it would bring the latter
into connection with the Bavarian triangles, and the triangle measurements
carried out in various parts of Europe, by being linked into a single whole, would, in a higher
\par\clearpage

% Printed 403
\noindent\null\hfill{\small\textbf{ON THE HANOVER TRIANGULATION.}}\hfill\makebox[0pt][r]{\fontsize{13}{16}\selectfont 403}\par\medskip
\noindent scientific context, become many times as valuable. That connection
between the Hanoverian and Electoral Hessian surveys was made in late 1823.
The latter have remained unfinished since then, although as much as was needed
for a bare connection of the Hanoverian with the Bavarian triangles, namely the
measurement of at least two angles in all the triangles needed for the connection,
was completed in 1823.

\par\noindent\hspace*{\parindent}Once at least a preliminary link between the northern and southern German
surveys had thus been achieved, it had to seem doubly important to establish
a link also with the great interconnected survey systems in the west. My
original commission was therefore extended to have me carry out a connection
from my arc-measurement triangles to \textsc{Krayenhoff}'s surveys as well. For
that purpose, in the years 1824 and 1825 I established a new system of triangles
from Hamburg to Jever. This established the connection with the Dutch, French,
and English triangles, so that all the great surveys extending across the most
cultivated part of all Europe are already in place.

\par\noindent\hspace*{\parindent}Since the Danish arc measurement had meanwhile remained unfinished,
I considered it necessary to arrange the astronomical part of my undertaking
so that the Hanoverian arc measurement could also stand on its own as a
complete whole.

\par\noindent\hspace*{\parindent}The Göttingen Observatory, itself a principal triangle station in the system
and the point from which all the triangle sides received their orientation,
naturally formed the southern endpoint. A circumstance unique of its kind
also made the choice of the northern endpoint a certainty: the observatory
that Professor \textsc{Schumacher} had meanwhile established in Altona lies almost
exactly on the meridian of the Göttingen Observatory. In 1827 I used the
\textsc{Ramsden} zenith sector at both places to observe 43 stars. I have already
published those observations, the difference in latitude derived from them, and
the other results in a work that appeared in 1828.

\par\noindent\hspace*{\parindent}To the trigonometric part of my undertaking the present work is
\par\smallskip
\noindent\hfill{\small 51\textsuperscript{*}}\hspace*{.122\linewidth}\par\clearpage

% Printed 404
\noindent\makebox[0pt][l]{\fontsize{13}{16}\selectfont 404}\hfill{\small\textbf{POSTHUMOUS PAPERS.}}\hfill\null\par\medskip
\noindent devoted. I considered it necessary for two reasons to give the account
all the detail it requires.

\par\noindent\hspace*{\parindent}In an \textls[120]{isolated} measurement of a meridian arc of moderate extent
\textls[120]{that is and is meant to be nothing more than that}, one can regard the
trigonometric part as subordinate to some degree. The accuracy attainable in
the astronomical observations falls far short of the accuracy attainable in
the trigonometric part, so it is less imperative to attain the highest degree
of accuracy in the latter.

\par\noindent\hspace*{\parindent}In a trigonometric survey of a country, on the other hand, it is advisable
for several reasons to determine the relative positions of the principal points
as accurately as the state of the art and the circumstances permit. In countless
cases this makes it possible to obtain sufficiently accurate derived determinations
of secondary points with extremely little work and by methods that would
otherwise yield wildly inaccurate results. When such a trigonometric survey
stands in isolation, a fuller publication of its constituent parts is admittedly
of no general interest. But the more the trigonometric surveys carried out in
different parts of Europe are connected with each other and gradually approach
one great whole, the more the individual parts take on the character of a
precious common possession of lasting value for all time, and the more
important it becomes to preserve all their essential details so fully that their
reliability can always be tested both as a whole and in detail.

\par\noindent\hspace*{\parindent}A second reason for going into detail was the distinctive character
of the methods I used, both in the measurements themselves and in processing
them to obtain results. Some of these methods differ entirely from those
usually employed, and describing them was to be a principal purpose of this
work. Combining a general account of this method with its continuous application
to an extensive survey system is undoubtedly the best means of showing
its nature in its true light and making it easier for those who wish to use it
in similar surveys in the future to do so.
\par\clearpage

% Printed 405
\noindent\null\hfill{\small\textbf{ON THE HANOVER TRIANGULATION.}}\hfill\makebox[0pt][r]{\fontsize{13}{16}\selectfont 405}\par\medskip
\begin{center}
\textls[120]{First Section.}\par\smallskip
General Arrangement of the Surveys.\par\medskip
1.
\end{center}

\par\noindent\hspace*{\parindent}The southern and northern parts of the region from Göttingen to Hamburg
differ greatly in character. The former is mountainous, with the summits
mostly wooded to a greater or lesser degree. Most towns and villages are
situated so that their towers afford a wide view in at most one direction,
if at all. The northern part, by contrast, is flat and broken in many places by
woodland, which makes it very difficult and often quite impracticable to use
isolated low rises. In most of this stretch there are no towns or villages at all
with towers suitable as triangle stations.

\begin{center}
2.
\end{center}

\par\noindent\hspace*{\parindent}Under these circumstances it could be foreseen that we could count on
hardly any natural triangle stations. At most of the stations we would have
to build dedicated signal towers or consider other artificial means of making
them visible.

\par\noindent\hspace*{\parindent}A principal consideration in this regard is the size intended for the
individual triangles. Clearly, in a very large-scale survey, covering a whole
continent, for example, it would generally be most advantageous for the
accuracy of the whole to make the triangles as large as possible. The same
applies to a survey of smaller extent insofar as it is regarded as a constituent
part of such a larger system. This claim, however, holds only on the assumption
that the angles in the largest triangles can be measured with the same precision
as those in small triangles, or at least with no appreciably lower precision.
This condition certainly does not hold when church towers or artificial
signal towers are used, and the moments at which such objects at very great
distances permit the greatest precision and certainty in the measurements
are extremely rare.

\begin{center}
\rule{.24\linewidth}{.4pt}
\end{center}
\clearpage

% Printed 406
\noindent\makebox[0pt][l]{\fontsize{13}{16}\selectfont 406}\hfill{\small\textbf{POSTHUMOUS PAPERS.}}\hfill\null\par
\vspace{2\baselineskip}
\begin{center}
[2.]\par\smallskip
{\large [Extracts from Reports on the Triangulation to the Hanoverian\\
Cabinet Ministry.]}
\end{center}
\bigskip
\begin{center}
[From a report dated 7~January 1822 »über die Arbeiten im Jahre 1821«.]
\end{center}

\par\noindent\hspace*{\parindent}Of the many operations that go into an arc measurement, forming the
network of triangles and measuring the angles is the one that requires by far
the most time and work. For the Hanoverian arc measurement this network
must connect with the Danish surveys in the north near Hamburg. In the south,
the Göttingen Observatory is indeed the true natural endpoint of the arc measurement
itself, but to make it possible to extend it farther south, it is also essential
to connect it with the foreign surveys that touch the Kingdom of Hanover
on the southern side. These include the surveys commissioned by the Royal
Prussian Government and carried out with great care, and a large triangulation
is now also contemplated in Electoral Hesse, to be carried out with all attainable
accuracy.

\par\noindent\hspace*{\parindent}The Prussian surveys just mentioned were communicated to me last
winter as far as they had then progressed. I also had the opportunity to obtain
a portion of the measurements made by the French Colonel \textsc{Epailly}
in 1804 and the following years in Hanover, especially in its southern part.
Having these and some other resources at my disposal could make the initial
selection of the triangle stations easier. I also considered that some of the
small difficulties and complications inseparable from the beginning of such
operations could always be dealt with more quickly and easily near Göttingen.
These considerations led me to begin the triangulation \textls[120]{on the southern side}.

\par\noindent\hspace*{\parindent}As early as 1820 I had made suitable arrangements to procure the
necessary instruments\textperiodcentered{} here I mention only
\par\clearpage

% Printed 407
\noindent\null\hfill{\small\textbf{ON THE HANOVER TRIANGULATION.}}\hfill\makebox[0pt][r]{\fontsize{13}{16}\selectfont 407}\par\medskip
\noindent those directly concerned with the geodetic part of the arc measurement.
For the actual angle measurements I had ordered a twelve-inch theodolite
from \textsc{Reichenbach} (whose workshop his former foreman \textsc{Ertel} has now
taken over entirely). Its completion and delivery had been promised for the
spring of 1821. I already had a smaller theodolite made by the English instrument
maker \textsc{Troughton}, which Professor \textsc{Schumacher} had kindly obtained for me.

\par\noindent\hspace*{\parindent}The means of making the triangle stations visible at very great distances
required special provision. As it was my intention, and was of the utmost
importance, to choose the largest possible triangles, the character of the
region through which they were to run left no hope of using many church
towers as triangle stations. Specially constructed signal towers have hitherto
been the most common means used in such cases. In practice, however, there
are frequent cases in which even this means proves inadequate. Such signal
towers, like church towers, are always very difficult to see and much more
difficult to observe at great distances where they are not silhouetted against
the sky, and especially where they are seen against a nearby dark-colored
background\textsuperscript{*}). For these and other reasons, other observers have often
made their angle observations at night, some exclusively so, having the distant
triangle stations made visible by means of large \textsc{Argand} lamps with very
precisely parabolic reflectors. These night observations do, of course, present
other great difficulties and inconveniences in turn. Particularly with very
large triangles, a successful observation must usually be bought, as it were,
at the price of many fruitless \mbox{attempts.\enspace\mbox{.\ .\ .\ .}}\quad Although I
was therefore disinclined to use this method of observation exclusively,
especially since my physical strength would hardly have been equal to the
hardships of continually spending the night on mountains that are mostly
high and difficult to reach, I nevertheless had to,

\par\smallskip\noindent\rule{.24\linewidth}{.4pt}\par\smallskip
{\footnotesize
\noindent\hspace*{\parindent}\textsuperscript{*}) My own experience last summer repeatedly confirmed this.
For example, during my entire stay of more than four weeks on the Brocken, I was able to see the signal tower built
on the Hils, \(7\tfrac{1}{2}\) miles away, only once or twice for a few minutes, and the church towers of Hanover,
12 miles away, not a single time, although the direction was precisely known.\par}
\clearpage

% Printed 408
\noindent\makebox[0pt][l]{\fontsize{13}{16}\selectfont 408}\hfill{\small\textbf{POSTHUMOUS PAPERS.}}\hfill\null\par\medskip
\noindent since no other means was yet known, be prepared in advance to use
it at least in some individual cases. I had therefore provisionally ordered three
such lamps from \textsc{Repsold} in Hamburg and from \textsc{Körner} in Jena. I received
these lamps in May 1821, and their performance in the various trials conducted
with them at suitable distances did indeed meet my expectations.

\par\noindent\hspace*{\parindent}With all the great difficulties mentioned in forming large triangles in
mind from the experience of others, even before I had acquired experience
of my own, I looked for an entirely new means of overcoming them.
Theoretical investigations had convinced me that sunlight reflected by even
very small plane mirrors had sufficient power to be visible at the greatest
distances and could be observed far more easily and accurately than any towers
and signals, even more readily than several \textsc{Argand} lamps combined at night.
Putting this idea to use required a special apparatus or instrument to direct
the reflected sunlight continuously, with the utmost accuracy and certainty,
toward any given point, however distant. Although I had already worked out
the essentials of the design for this purpose in full, it was difficult to have
the instrument made, especially if this had to be done by an instrument maker
elsewhere. That would have deprived me of the advantage of ongoing oral
consultation on particular technical difficulties and would at least have delayed
completion considerably. This difficulty was, however, close to being resolved
when our skilled inspector \textsc{Rumpf}, who had been away for much of the winter,
returned to Göttingen around Easter and shortly afterward undertook to build
such an \mbox{instrument.\enspace\mbox{.\ .\ .\ .\ .\ .}}

\par\noindent\hspace*{\parindent}A particularly essential and important circumstance in the Hanoverian
arc measurement is that the first triangle sides running out from the observatory
here can be oriented with the greatest accuracy by means of the observatory's
own fixed meridian instruments. Originally, however, the observatory's view
in the direction of the meridian was clear neither to the north nor to the south.
Had this defect been impossible to remedy, that most important advantage would
have been entirely
\par\clearpage

% Printed 409
\noindent\null\hfill{\small\textbf{ON THE HANOVER TRIANGULATION.}}\hfill\makebox[0pt][r]{\fontsize{13}{16}\selectfont 409}\par\medskip
\noindent lost, and the defect would also have remained forever a fundamental
error in the choice of the observatory site. Fortunately, one half of this difficulty
had already been overcome. In the autumn of 1820 I had already opened up
the view to the north, which had previously been blocked by the gardens
outside Göttingen, and had a provisional meridian mark set up on a hill
near Weende. The difficulties on the south side of the meridian were far greater.
There a dense, tall forest three hours away limited the view. This obstacle
had to be overcome if \mbox{possible.\enspace\mbox{.\ .\ .\ .\ .\ .}}

\par\noindent\hspace*{\parindent}During this time Inspector \textsc{Rumpf} had already been working diligently
on the instrument mentioned above, which I shall henceforth call by the name
given to it, heliotrope. Before it was finished, however, I had hit on the idea
of adapting a simple reflecting sextant to serve as a kind of auxiliary heliotrope,
admittedly far less perfect than that instrument itself, yet also usable when
handled skillfully. Trials already conducted with it at distances of nearly
2 miles confirmed the enormous power of reflected sunlight to a degree that
almost exceeded my \mbox{expectations.\enspace\mbox{.\ .\ .\ .\ .\ .}}
\par\bigskip

\begin{center}
[From a report dated 31~January 1823 »über die Arbeiten der Gradmessung im Jahre 1822«.]
\end{center}

\par\noindent\hspace*{\parindent}In 1821 the triangulation, the most extensive part of the whole undertaking,
was begun, and the measurements at five triangle stations were completed,
namely at the Göttingen Observatory, at the meridian mark, on the Hohehagen,
the Hils, and the Brocken. To continue the work, four further points farther
north were selected, namely Lichtenberg, the Deister, the Wohlenberg in the
district of Gifhorn, and a hill near Brelingen in the bailiwick of Bissendorf.
I myself had not yet visited these points, however, and it remained doubtful
whether some of them would qualify as triangle stations. That had to be settled
before a plan for the first work of 1822 could be drawn up.

\par\smallskip
\noindent\hspace*{.030\linewidth}{\footnotesize IX.}\hfill{\small 52}\hspace*{.137\linewidth}\par
\clearpage

% Printed 410
\noindent\makebox[0pt][l]{\fontsize{13}{16}\selectfont 410}\hfill{\small\textbf{POSTHUMOUS PAPERS.}}\hfill\null\par\medskip
\par\noindent\hspace*{\parindent}But this was only the smallest part of the necessary preliminary
investigations. The operations were now approaching the Lüneburg Heath,
an entirely flat region where the lack of commanding points and the almost
countless larger and smaller woods covering it like a chessboard make the
formation of triangles of any substantial size extraordinarily difficult. I already
knew of these difficulties from the reports of the French Colonel \textsc{Epailly},
who had directed the French surveys in the Electorate of Hanover in 1804
and 1805. This engineer had judged the difficulties of the terrain so great
that he declared it impossible to form a triangle system from the Aller to
the Elbe. He therefore made the connection by an enormous detour, namely
by triangles along the Weser as far as its mouth and then back up the Elbe
to Hamburg. This procedure might be acceptable at most as a makeshift in
a land survey and could never be admissible in an arc measurement.

\par\noindent\hspace*{\parindent}These circumstances made a reconnaissance journey necessary before
the actual surveying operations began. This journey, on which Captain
\textsc{Müller} was the only one of my three assistants I took with me, occupied me
from 28~April to 1~June. I pass over the details of the laborious investigations
here and give only their principal results.

\par\noindent\hspace*{\parindent}The Brelingerberg was found unsuitable for making the connection. This
was more than made up for by the fortunate discovery that two new points
farther north, a high field near Garssen and the Falkenberg, the former almost
one mile northeast of Celle and the latter one mile northwest of Bergen, could
both be connected directly with Lichtenberg and with the Deister. This made
the Wohlenberg superfluous, and the initial surveying operations were now definitely and
securely determined. As for proceeding farther north, I indeed found the difficulties
as great as I had expected, yet I also managed to establish, as it were in the
heart of the heath, the feasibility of two good triangles between the four
points Falkenberg, Hauselberg (in the bailiwick of Hermannsburg), Wulfsode
(in the district of Ebstorf), and Wilsede (on the far southwestern boundary of
the district of Winsen an der Luhe). It
